\documentclass[noauthor,noinstructornotes]{ximera}

\title{Week 3}
\author{Math 1193 Design Team}

\begin{document}
\begin{abstract}
    Activities for Week 3.
\end{abstract}
\maketitle

\begin{problem}
Solve \(4x - 5y = 10\) for \(x\). Now solve it for \(y\).
\end{problem}

\begin{problem}
Solve \(0.01x + 0.5y = 1\) for \(x\). Now solve it for \(y\).
Share with a group how you solved the equation.
\end{problem}

\begin{problem}
\begin{enumerate}
    \item Graph the equations \(x + y = 2\) and \(x + y = 3\)
    without a calculator. Do these lines ever touch? Why or why not?
    Can you find an \(x\) and a \(y\) that satisfy both equations at 
    the same time? Why or why not? 
    \item Graph the equations \(x + y = 2\) and \(x - y = 3\)
    without a calculator. Do these lines ever touch? Why or why not?
    Can you find an \(x\) and a \(y\) that satisfy both equations at 
    the same time? Why or why not? 
    \item Is it possible for two lines in the 
    plane to touch more than once? Debate with a friend.
\end{enumerate}

\begin{instructorNotes}
Hopefully this brings out different methods students use to graph lines.
\end{instructorNotes}
\end{problem}

\begin{problem}
A theater charges \$50 per ticket for seats in Section A, 
\$30 per ticket for seats in Section B, and \$20 per ticket 
for seats in Section C. For one play, 4000 tickets were sold 
for a total of \$120,000 in revenue. 

\begin{enumerate}
    \item If 1000 more tickets in Section B were sold 
    than the other two tickets combined, how many tickets 
    in each section were sold?
    \item Define variables that help you solve this 
    problem. Be specific about what quantities the variables
    represent. Compare with a friend and discuss why you made 
    the choices you made.
    \item Write equations that relate the variables you 
    have defined.
\end{enumerate}

\begin{instructorNotes}
Closely mirrors 1148 exam-type questions with a bit more scaffolding.
\end{instructorNotes}
\end{problem}

\end{document}
